% §11.2 Las Vegas vs Monte Carlo
%
% Distinguishes the two species of randomized algorithm by where the
% randomness is allowed to land: in the runtime (Las Vegas) or in the output
% (Monte Carlo). Frames the trade-off, surveys three canonical examples
% (randomized quicksort, Karger min-cut, Miller-Rabin), discusses
% conversions in both directions, and positions Monte Carlo decision
% algorithms as $\Bernoulli{\BoolType}$ instances within the framework of
% Part~II.
%
% Cross-reference resolution status:
%
% RESOLVED (chapters 3, 4, drafted in Plans 1-4):
%   \Cref{def:bernoulli-bool}       -- chapter 4 §4.5 $\Bernoulli{\BoolType}$
%   \Cref{thm:composition}          -- chapter 3 §3.3 composition theorem
%   \Cref{sec:miller-rabin}         -- chapter 3 §3.4 informal walkthrough
%
% INTERNAL (chapter 11):
%   \Cref{sec:miller-rabin-formal}  -- §11.1 (Plan 11 Task 2, just drafted)
%   \Cref{thm:miller-rabin-formal}  -- §11.1 theorem
%   \Cref{sec:rp-bpp}               -- §11.3 (Plan 11 Task 4, forward)
%   \Cref{sec:sample-complexity}    -- §11.4 (Plan 11 Task 5, forward)

\section{Las Vegas vs Monte Carlo}
\label{sec:las-vegas-vs-monte-carlo}

The chapter~3 \S3.4 walkthrough and the formalization of
\Cref{sec:miller-rabin-formal} both treat a single example: Miller-Rabin.
That example happens to live at one corner of the randomized-algorithm
landscape. To navigate the rest of the landscape one needs a coordinate
system, and the standard one separates randomized algorithms by where the
randomness is allowed to surface. It either surfaces in the runtime, while
the output stays exactly correct, or it surfaces in the output, while the
runtime stays predictably bounded. These two regimes are named after the
casinos that gave them their public reputations: Las Vegas and Monte Carlo.
This section lays out the two definitions, shows the trade-off that forces
the choice, walks through three canonical examples, sketches the
conversions between the two regimes, and identifies which regime maps onto
the $\Bernoulli{\BoolType}$ vocabulary that the rest of Part~III speaks.

%---------------------------------------------------------------------------
\subsection*{Two definitions}

\begin{definition}[Las Vegas algorithm]
\label{def:las-vegas}
A randomized algorithm $A$ on inputs $x$ is a \emph{Las Vegas algorithm}
when its output $A(x)$ is deterministically correct on every input and
every coin sequence (whenever $A$ halts), and its runtime $T_A(x)$ is a
random variable whose distribution may depend on $x$. The algorithm is
\emph{expected-polynomial-time Las Vegas} when $\Expect{T_A(x)}$ is bounded
by a polynomial in $|x|$ for every $x$.
\end{definition}

The contract a Las Vegas algorithm signs is asymmetric. Output correctness
is absolute: no coin flip will ever cause a wrong answer. Runtime is a
random variable, and a bad coin sequence can make the algorithm take longer
than its expected value. There is no fixed worst-case time bound; an
adversarial coin flip is allowed to extend the runtime, as long as the
expected runtime stays inside the promised polynomial envelope.

\begin{definition}[Monte Carlo algorithm]
\label{def:monte-carlo}
A randomized algorithm $A$ on inputs $x$ is a \emph{Monte Carlo algorithm}
when its runtime $T_A(x)$ is deterministically bounded on every input and
every coin sequence (typically by a polynomial in $|x|$), and its output
$A(x)$ is a random variable that agrees with the true answer $f(x)$ only
with bounded probability. The probability of disagreement is the
algorithm's \emph{error rate}, and Monte Carlo algorithms are classified
further by whether that rate is one-sided or two-sided. A one-sided Monte
Carlo decision algorithm has $\fprate \le \alpha$ and $\fnrate = 0$ (or the
mirror image); a two-sided one has $\fprate, \fnrate \le \alpha$ for some
common bound $\alpha < 1/2$.
\end{definition}

The Monte Carlo contract is the mirror image. Runtime is locked down: every
coin sequence completes in at most $T(|x|)$ steps. Output is random, and a
bad coin sequence is allowed to return the wrong answer, as long as the
probability of that bad sequence is below the advertised error rate.

The two definitions are exhaustive of the practically useful regimes but
not of the formal possibilities. A randomized algorithm whose runtime is
unbounded \emph{and} whose output may be wrong is rarely interesting; it
loses both contracts simultaneously and offers nothing in exchange. The
opposite, an algorithm with bounded worst-case runtime \emph{and}
deterministically correct output, is simply a deterministic algorithm and
needs no randomness at all. Las Vegas and Monte Carlo carve up the
non-trivial middle.

%---------------------------------------------------------------------------
\subsection*{The trade-off}

The two regimes look like a matter of taste but they are not. They are
two faces of a single trade-off, which the following observation makes
precise.

\begin{remark}
If $f\colon X \to Y$ is a total function whose graph an algorithm wishes to
compute, then for any randomized algorithm $A$ at most one of the following
holds with probability $1$ on every input: (i) $A(x) = f(x)$, or (ii)
$T_A(x) \le T(|x|)$ for a fixed function $T$. An algorithm that satisfies
both for every input on every coin sequence is deterministic.
\end{remark}

The remark is almost a tautology, but its content is the design tension of
the field. A worst-case bounded runtime over all coin sequences leaves the
algorithm no room to keep flipping coins until the answer is right. If
every coin sequence completes in at most $T(|x|)$ steps, and on some coin
sequence the answer the algorithm has computed at step $T(|x|)$ happens to
be wrong, the algorithm must return that wrong answer. Conversely, an
algorithm that insists on returning the right answer must keep working
until it is sure, and on bad coin sequences that takes longer than any
fixed bound would allow.

Each regime gives up exactly one of the two guarantees. Las Vegas keeps the
output correct and pays in the runtime tail: a bad sample of pivots, a
bad sample of hash seeds, and the algorithm runs longer than expected.
Monte Carlo keeps the runtime bounded and pays in the output: a bad sample
of witnesses, a bad sample of contractions, and the algorithm reports a
wrong answer. Which side to pay on is a question about the application.
When an answer that is sometimes wrong is unsafe (a medical decision, a
financial settlement) but a longer wait is tolerable, Las Vegas is the
sensible choice. When timing is part of the contract and a small error
rate is tolerable (a network filter, a real-time approximate query), Monte
Carlo is.

%---------------------------------------------------------------------------
\subsection*{Three examples}

The two definitions become concrete once one has a few canonical examples
in mind. The following three are representative and recur throughout the
rest of Part~III.

\paragraph{Randomized quicksort (Las Vegas).} Sort an array of $n$
elements by recursively partitioning around a pivot chosen uniformly at
random from the current sub-array. The output is always the correctly
sorted array: every comparison was a comparison, and the partition step
preserves the multiset of elements. Runtime is a random variable. The
expected number of comparisons is $\Expect{T(n)} = O(n \log n)$, but a
bad sequence of pivot choices can make a particular run take $\Theta(n^2)$
comparisons. The worst-case sequence is exponentially unlikely:
$\Pr[T(n) \ge c n \log n] \le n^{-\Omega(c)}$ for a suitable constant. The
algorithm is Las Vegas because the contract is on output correctness, and
the runtime tail is the price paid.

\paragraph{Karger min-cut (Monte Carlo).} Find a global minimum cut of an
$n$-vertex multigraph by repeatedly contracting a uniformly random edge
until only two super-vertices remain, then reporting the set of edges
between them as the cut. The runtime is deterministically $O(n^2)$ on
every coin sequence. The output is random; on a fixed input graph $G$ with
a fixed minimum cut $C$, Karger's analysis gives
\[
  \Pr[\text{algorithm returns } C] \;\ge\; \frac{2}{n(n-1)}.
\]
A single run has only a $\Theta(1/n^2)$ chance of returning the right cut.
The algorithm is salvaged by amplification: running it $k = c n^2 \log n$
independent times and taking the smallest cut found drives the failure
probability below $1/n^c$. Each run is fixed-cost; the cost of confidence
is the count of runs. This is Monte Carlo through and through.

\paragraph{Miller-Rabin (Monte Carlo).} The primality test of
\Cref{sec:miller-rabin-formal}. The $k$-round version has runtime
deterministically bounded by $O(k \log^3 n)$ on every coin sequence (each
round does a modular exponentiation, which is $O(\log^3 n)$ with
schoolbook modular arithmetic). The output is random: on composite inputs,
the algorithm may report \textsc{probably-prime} with probability at most
$(1/4)^k$, by \Cref{thm:miller-rabin-formal}. On prime inputs, the
algorithm always reports \textsc{probably-prime}, so the error is
one-sided. This is the Monte Carlo regime, and it is exactly the regime in
which the $\Bernoulli{\BoolType}$ vocabulary of Part~II applies directly.

The three examples cover the corners of the design space. Randomized
quicksort sacrifices runtime tail to keep output exact. Karger min-cut
sacrifices output to keep runtime fixed and uses amplification to recover
confidence. Miller-Rabin sacrifices output in a one-sided way and uses the
composition theorem of chapter~3 to amplify confidence in the same way.

%---------------------------------------------------------------------------
\subsection*{Conversion in both directions}

The Las Vegas and Monte Carlo regimes are not strictly disjoint as a matter
of design. Each can be converted to the other, at a cost.

\paragraph{Las Vegas to Monte Carlo.} Given a Las Vegas algorithm $A$ with
expected runtime $\Expect{T_A(x)} \le \mu(|x|)$, define the Monte Carlo
algorithm $A'$ that runs $A$ for at most $c \cdot \mu(|x|)$ steps and
returns ``failure'' if $A$ has not halted by then. By Markov's inequality,
$A'$ returns failure with probability at most $1/c$, and otherwise returns
the same correct output as $A$. Picking $c = 10$ gives a Monte Carlo
algorithm with error rate $\le 0.1$ and worst-case runtime $\le 10\,
\mu(|x|)$. The conversion is mechanical, and it explains why expected
polynomial time is essentially the same resource as worst-case polynomial
time with a bounded error rate.

\paragraph{Monte Carlo to Las Vegas (if a verifier exists).} Given a Monte
Carlo algorithm $A$ with error rate $\fprate$ \emph{and} a deterministic
polynomial-time verifier $V$ that, given an input $x$ and a candidate
output $y$, decides whether $y = f(x)$, define the Las Vegas algorithm
$A''$ that repeatedly runs $A(x)$, verifies the output with $V$, and
returns the first verified result. The expected number of trials is
$1/(1-\fprate)$, so the expected runtime is bounded. The output is always
correct because $V$ accepts only the true answer. The conversion requires
the verifier, which is a non-trivial assumption: some problems admit Monte
Carlo algorithms but no known polynomial-time verifier. This asymmetry is
why the complexity class ZPP (the Las Vegas analogue of P) is contained in
both RP and coRP, but the reverse containment is open in general; see
\Cref{sec:rp-bpp}.

%---------------------------------------------------------------------------
\subsection*{The $\Bernoulli{\BoolType}$ correspondence}

The Monte Carlo regime is the one that Part~II's vocabulary fits exactly.
A Monte Carlo \emph{decision} algorithm, on input $x$, returns a Boolean
report $A(x) \in \{\top, \bot\}$ that is the algorithm's guess at the true
predicate value $P(x) \in \{\top, \bot\}$. The error structure is precisely
the $2 \times 2$ channel of \Cref{def:bernoulli-bool}:
\[
  \fprate \;=\; \Pr[A(x) = \top \mid P(x) = \bot],
  \qquad
  \fnrate \;=\; \Pr[A(x) = \bot \mid P(x) = \top].
\]
A Monte Carlo decision algorithm \emph{is} a $\Bernoulli{\BoolType}$
instance, in the sense of the type-membership criterion of chapter~4. The
algorithm's stated error rate is the framework's $\fprate$ (or its mirror
$\fnrate$ in the symmetric case). The $k$-fold amplification calculation
of \Cref{thm:miller-rabin-formal} is just the composition theorem applied
to $k$ independent instances of the same Bernoulli predicate. Everything
chapter~7 says about PPV, NPV, and the derived measures applies to a Monte
Carlo decision algorithm without modification, with prevalence playing the
role of the base rate at which the input distribution presents true
positives.

Las Vegas algorithms occupy a different level of abstraction. The output
of a Las Vegas algorithm is not a Bernoulli random variable; it is
deterministically equal to the true answer. The randomness lives in the
runtime, which is a non-negative real-valued random variable, not a
Boolean. Framework-level objects such as $\fprate$ and $\fnrate$ are not
defined for a Las Vegas algorithm because there is no error to rate. The
right framework for analyzing a Las Vegas algorithm is the theory of
runtime tail bounds: Markov's inequality, Chebyshev's inequality, and the
Chernoff-style concentration inequalities we will see in
\Cref{sec:sample-complexity}. These are about the runtime distribution,
not about output error.

The asymmetry is principled. Part~III is about the algebra of approximate
answers, and approximate answers are what Monte Carlo algorithms produce.
Las Vegas algorithms produce exact answers and then ask the user to wait;
their analysis belongs to a different chapter of algorithmics, the one
about expected runtime, and the Bernoulli framework has nothing to add to
it. From this point on in Part~III, when we say ``randomized algorithm''
without further qualification, we mean a Monte Carlo decision algorithm,
and we mean it as a $\Bernoulli{\BoolType}$ instance.
